\subsection{Class CodeSource}
\label{sec:codesource}

\declaredin{CodeSource.h}

The CodeSource interface is used by the ParseAPI to retrieve binary code from
an executable, library, or other binary code object; it also can provide hints
of function entry points (such as those derived from debugging symbols) to seed
the parser. The ParseAPI provides a default implementation based on the
SymtabAPI that supports many common binary formats. For details on implementing
a custom CodeSource, see Appendix~\ref{sec:extend}.

\begin{apient}
virtual bool nonReturning(Address func_entry);
bool nonReturning(std::string func_name);
\end{apient}
\apidesc{Looks up whether a function returns (by name or location). This information may be statically known for some code sources, and can lead to better parsing accuracy.}

\begin{apient}
virtual bool nonReturningSyscall(int number);
\end{apient}
\apidesc{Looks up whether a system call returns (by system call number). This information may be statically known for some code sources, and can lead to better parsing accuracy.}


\begin{apient}
virtual Address baseAddress() const;
virtual Address loadAddress() const;
\end{apient}
\apidesc{If the binary file type supplies non-zero base or load addresses (e.g. Windows PE), implementations should override these functions.}

\begin{apient}
const std::map<Address, std::string>& linkage() const;
\end{apient}
\apidesc{Returns a reference to the external linkage map, which may or may not be filled in for a particular CodeSource implementation.}

\begin{apient}
struct Hint {
    Address _addr;
    int _size;
    CodeRegion* _reg;
    std::string _name;

    Hint();
    Hint(Address a, int size, CodeRegion* r, std::string s);
    bool operator<(const Hint& h) const;
};

dyn_c_vector<Hint> const& hints() const;
\end{apient}
\apidesc{Returns a vector of the currently defined function entry hints. A hint
names a function entry at \code{\_addr} within region \code{\_reg}, with
\code{\_size} the size of the function taken from its symbol, or 0 where that
is not known. Hints compare by address.
\code{dyn\_c\_vector} is a oneTBB concurrent vector
{\raggedright\url{https://www.intel.com/content/www/us/en/docs/onetbb/developer-guide-api-reference/2021-9/concurrent-vector.html}.\par}}

\begin{apient}
std::vector<CodeRegion*> const& regions() const;
\end{apient}
\apidesc{Returns a read-only vector of code regions within the binary represented by this code source.}

\begin{apient}
int findRegions(Address addr,
                std::set<CodeRegion*>& ret) const;
\end{apient}
\apidesc{Finds all CodeRegion objects that overlap the provided address. Some code sources (e.g. archive files) may have several regions with overlapping address ranges; others (e.g. ELF binaries) do not.}

\begin{apient}
bool regionsOverlap() const;
\end{apient}
\apidesc{Indicates whether the CodeSource contains overlapping regions.}
